\chapter{Neural Networks for Noisy Tabular Data}\label{ch:ml:neural-networks-for-noisy-tabular-data}

Two researchers train the same small network, on the same data, with the same code, and report information coefficients
of 0.057 and 0.065 on the same ten test years. They are arguing about whether a third hidden layer helps; on average it
adds 0.002. The only difference between their two runs is the random seed, and it is worth four times the architecture
they are arguing about. On market data a \omterm{def:ml:neural-networks-for-noisy-tabular-data:mlp}{neural network} is a random variable before it is a model: its initial weights,
its \omterm{def:ml:neural-networks-for-noisy-tabular-data:dropout}{dropout} masks and the order of its \omterm{def:ml:neural-networks-for-noisy-tabular-data:backprop}{mini-batches} all move the result, and the signal is too weak to drown them out.
This chapter builds the \omterm{def:ml:neural-networks-for-noisy-tabular-data:mlp}{multilayer perceptron}, trains it the way noisy tabular data require, and measures its seeds as a
source of variance.

\section{The multilayer perceptron and backpropagation}

\begin{definition}[Neural network, multilayer perceptron, activation function]\label{def:ml:neural-networks-for-noisy-tabular-data:mlp}
A \emph{neural network}\index{neural network} is a function built by composing parametrised linear maps with fixed
nonlinear functions. A \emph{multilayer perceptron}\index{multilayer perceptron} (MLP) is the fully connected case:
$z_0 = x$, $z_l = \mathrm{act}(A_lz_{l-1} + c_l)$ for $l = 1,\dots,L$, and $f_\theta(x) = A_{L+1}z_L + c_{L+1}$, with
parameters $\theta = (A_l, c_l)_l$. The \emph{activation function}\index{activation function} $\mathrm{act}$ is applied
coordinate by coordinate; the usual one is $\mathrm{relu}(u) = \max(u, 0)$.
\end{definition}

\begin{omfigure}
\begin{tikzpicture}[x=2.4cm,y=0.55cm,font=\footnotesize,
  neuron/.style={circle,draw=omDef,fill=omDef!12,minimum size=9pt,inner sep=0pt}]
\foreach \i in {1,...,6} { \node[neuron] (a\i) at (0,\i-3.5) {}; }
\foreach \i in {1,...,5} { \node[neuron] (b\i) at (1,\i-3) {}; }
\foreach \i in {1,...,3} { \node[neuron] (c\i) at (2,\i-2) {}; }
\node[neuron,draw=omThm,fill=omThm!15] (d) at (3,0) {};
\foreach \i in {1,...,6} { \foreach \j in {1,...,5} { \draw[gray!45] (a\i) -- (b\j); } }
\foreach \i in {1,...,5} { \foreach \j in {1,...,3} { \draw[gray!45] (b\i) -- (c\j); } }
\foreach \i in {1,...,3} { \draw[gray!45] (c\i) -- (d); }
\node[align=center] at (0,-4.6) {40 ranked\\characteristics};
\node[align=center] at (1,-4.6) {32 units\\$\mathrm{relu}(A_1x + c_1)$};
\node[align=center] at (2,-4.6) {16 units\\$\mathrm{relu}(A_2z_1 + c_2)$};
\node[align=center] at (3,-4.6) {forecast\\$A_3z_2 + c_3$};
\end{tikzpicture}

\omcaption{The chapter's small network (drawn with fewer units than it has): 40 inputs, hidden layers of 32 and 16
units, one output, about 1\,850 parameters. The larger network adds a first layer of 64 units.}\label{fig:ml:nets:mlp}
\end{omfigure}

\begin{definition}[Backpropagation, mini-batch, epoch]\label{def:ml:neural-networks-for-noisy-tabular-data:backprop}
\emph{Backpropagation}\index{backpropagation} computes the gradient of the loss with respect to every parameter of a
network by reverse-mode algorithmic differentiation (Book~4, chapter~28): one forward pass stores the layers' values, one
backward pass propagates $\partial\mathcal L/\partial z_l$ from the output to the input. A \emph{mini-batch}\index{mini-batch}
is the random subset of training rows on which one gradient step is computed; an \emph{epoch}\index{epoch} is one pass
over the \omterm{def:ml:why-financial-machine-learning-is-different:sets}{training set} in mini-batches.
\end{definition}

The cost of the gradient is a small multiple of the cost of the forward pass, whatever the number of parameters: that is
Book~4's result on reverse mode, and it is why networks can have millions of parameters. The chapter's networks have two and
five thousand and fit in about a second on one CPU thread.

\section{Optimisers, normalisation and regularisation}

\begin{definition}[Adam, weight decay]\label{def:ml:neural-networks-for-noisy-tabular-data:adam}
\emph{Adam}\index{Adam} is stochastic gradient descent (Book~4, chapter~24) with a per-parameter step: with gradient
$g_k$ at step $k$, it keeps averages $m_k = \beta_1m_{k-1} + (1-\beta_1)g_k$ and $v_k = \beta_2v_{k-1} + (1-\beta_2)g_k^2$
and moves $\theta$ by $-\eta\,\hat m_k/(\sqrt{\hat v_k} + \epsilon)$, where hats undo the averages' bias towards zero
(Kingma and Ba). \emph{Weight decay}\index{weight decay} shrinks every weight by $\eta\lambda\theta$ at each step, apart
from the gradient step (AdamW, Loshchilov and Hutter); for plain gradient descent it is a ridge penalty.
\end{definition}

\begin{definition}[Dropout, batch and layer normalisation]\label{def:ml:neural-networks-for-noisy-tabular-data:dropout}
\emph{Dropout}\index{dropout} sets each hidden unit to zero with probability $d$ during training, scaling the others by
$1/(1 - d)$, and uses the full network at prediction. \emph{Batch normalisation}\index{batch normalisation} standardises
each hidden unit over the rows of the \omterm{def:ml:neural-networks-for-noisy-tabular-data:backprop}{mini-batch} (running averages at prediction); \emph{layer
normalisation}\index{layer normalisation} standardises the units of each row, independently of the batch.
\end{definition}

The chapter's data are chapter~4's panel (500 stocks, forty ranked characteristics with factor risk, a ceiling of
0.71\%): the networks are trained on months 0 to 199 with \omterm{def:ml:trees-and-boosting:early}{early stopping} on months 201 to 239 (chapter~5's rule, one
month purged), and scored on the ten test years with \texttt{firm.mlbase}'s report, exactly as chapter~4's models were.
\omterm{def:ml:trees-and-boosting:early}{Early stopping} stops them after two to four \omterm{def:ml:neural-networks-for-noisy-tabular-data:backprop}{epochs} (\cref{fig:ml:nets:history}): on data this noisy, the validation
score peaks almost at once, and every later \omterm{def:ml:neural-networks-for-noisy-tabular-data:backprop}{epoch} fits more noise than signal.

\begin{omfigure}
\begin{tikzpicture}
\begin{axis}[width=10.5cm,height=5.2cm,xlabel={epoch},ylabel={validation $R^2$ (\%)},tick label style={font=\small},
  label style={font=\small},xmin=0.5,xmax=9.5,ymin=0.3,ymax=0.85,xtick={1,2,3,4,5,6,7,8,9},grid=major,
  grid style={gray!25},legend columns=3,legend style={font=\footnotesize,at={(0.5,-0.3)},anchor=north,draw=none,
  fill=none,/tikz/every even column/.append style={column sep=8pt}},table/col sep=comma,unbounded coords=jump]
\addplot[omDef,very thick,mark=*] table[x=epoch,y=s1]{figdata/ml/07-neural-networks-for-noisy-tabular-data/history.csv};
\addplot[omThm,very thick,mark=square*] table[x=epoch,y=s2]{figdata/ml/07-neural-networks-for-noisy-tabular-data/history.csv};
\addplot[omProp,very thick,mark=triangle*] table[x=epoch,y=s3]{figdata/ml/07-neural-networks-for-noisy-tabular-data/history.csv};
\legend{seed 1,seed 2,seed 3}
\end{axis}
\end{tikzpicture}

\omcaption{Validation R-squared of the small network after each \omterm{def:ml:neural-networks-for-noisy-tabular-data:backprop}{epoch}, three seeds; training stops five \omterm{def:ml:neural-networks-for-noisy-tabular-data:backprop}{epochs} after the
best. The seeds peak at \omterm{def:ml:neural-networks-for-noisy-tabular-data:backprop}{epochs} 2, 2 and 4. Data: \texttt{ml\_nets.fitted}.}\label{fig:ml:nets:history}
\end{omfigure}

With three seeds each, \omterm{def:ml:neural-networks-for-noisy-tabular-data:dropout}{dropout} of 0.2 raises the mean test IC from 0.059 to 0.062, \omterm{def:ml:neural-networks-for-noisy-tabular-data:adam}{weight decay} of 0.01 changes nothing
at this scale, and both normalisations hurt: \omterm{def:ml:neural-networks-for-noisy-tabular-data:dropout}{batch normalisation} to 0.052, \omterm{def:ml:neural-networks-for-noisy-tabular-data:dropout}{layer normalisation} to 0.056. Normalising
hidden units helps deep networks train; these are shallow, their inputs are already ranks, and a batch's statistics are
one more source of noise. The differences are a few thousandths of IC, comparable to the seed noise measured next, and
three seeds per variant is too few to rank them with confidence.

\section{Seeds as a source of variance}

Ten seeds of the small network (32 and 16 units) score test ICs from 0.057 to 0.065, with a mean of 0.0601 and a
standard deviation of 0.0023; ten seeds of the larger one (64, 32 and 16) average 0.0619 (\cref{fig:ml:nets:seeds}).
Their R-squared varies far more than their IC, from 0.08\% to 0.30\% for the small network: the seed moves the scale of
the forecasts, which R-squared punishes and a ranking ignores (chapter~1). The decile portfolio's net Sharpe ratio goes
from 1.37 to 1.78.

\begin{proposition}[How many seeds a comparison needs]\label{prop:ml:nets:seeds}
If two architectures' scores vary across seeds with standard deviation $s$, independently, and their means differ by
$\delta$, the difference of their $n$-seed averages has standard error $s\sqrt{2/n}$, so it reaches two standard errors
when $n \ge 8s^2/\delta^2$.
\end{proposition}

\begin{proof}
$\Var(\bar a - \bar b) = s^2/n + s^2/n$; solve $\delta = 2s\sqrt{2/n}$ for $n$.
\end{proof}

With $s = 0.0024$ (pooled) and $\delta = 0.0019$, the comparison needs 14 seeds per architecture. A single run of each,
the usual practice, compares two draws whose difference has a standard error of 0.0035, almost twice the effect.

\begin{definition}[Deep ensemble]\label{def:ml:neural-networks-for-noisy-tabular-data:ensemble}
A \emph{deep ensemble}\index{deep ensemble} averages the predictions of networks of the same architecture trained from
different random seeds on the same data (Lakshminarayanan, Pritzel and Blundell, 2017).
\end{definition}

Averaging removes the variance the seeds put in. The ten-seed ensemble of the small network scores an IC of 0.066, above
every one of its members, with an R-squared of 0.32\% and a net Sharpe ratio of 1.73; the ensemble of the larger
network scores 0.066 as well. The architecture question the two researchers argued about disappears once each network
is averaged over its seeds: the averaging bought what the architecture could not. It is \cref{prop:ml:trees:bagging}
again, with seeds in the role of bootstrap samples.

\begin{omfigure}
\begin{tikzpicture}
\begin{axis}[width=9cm,height=5.6cm,ylabel={test information coefficient},xmin=0.4,xmax=2.6,ymin=0.054,ymax=0.068,
  xtick={1,2},xticklabels={{32, 16 units},{64, 32, 16 units}},tick label style={font=\small},label style={font=\small},
  grid=major,grid style={gray!25},yticklabel style={/pgf/number format/fixed,/pgf/number format/precision=3},
  scaled y ticks=false,legend columns=3,legend style={font=\footnotesize,at={(0.5,-0.2)},anchor=north,draw=none,
  fill=none,/tikz/every even column/.append style={column sep=8pt}},table/col sep=comma]
\addplot[only marks,mark=*,mark size=1.8pt,gray] table[x=arch,y=ic]{figdata/ml/07-neural-networks-for-noisy-tabular-data/seeds.csv};
\addplot[only marks,mark=-,mark size=9pt,very thick,omDef] table[x=arch,y=mean]{figdata/ml/07-neural-networks-for-noisy-tabular-data/means.csv};
\addplot[only marks,mark=star,mark size=4pt,thick,omThm] table[x=arch,y=ensemble]{figdata/ml/07-neural-networks-for-noisy-tabular-data/means.csv};
\legend{one seed,mean of 10 seeds,ensemble of 10}
\end{axis}
\end{tikzpicture}

\omcaption{Test IC of ten seeds of two architectures on chapter~4's panel, their means, and the IC of each ten-seed
ensemble. Data: \texttt{ml\_nets.seed\_study} and \texttt{ml\_nets.ensemble}.}\label{fig:ml:nets:seeds}
\end{omfigure}

\section{Ensembles}

A seed ensemble costs $n$ trainings and nothing else, and it is embarrassingly parallel. It also gives something a
single network cannot: the spread of its members' forecasts for each row, which chapter~11 uses as a measure of the
model's own uncertainty. Ensembles of different architectures, or of networks with boosted trees and a ridge
regression (stacking, Book~7, chapter~14), go further when the members' errors are less correlated than seeds' are.

\begin{method}[Training a network on noisy tabular data]\label{met:ml:nets:training}
\begin{enumerate}
\item Rank the inputs by date; scale the target by its training standard deviation.
\item Start small (two hidden layers, tens of units), with AdamW, a \omterm{def:ml:trees-and-boosting:boosting}{learning rate} near $10^{-3}$, and \omterm{def:ml:trees-and-boosting:early}{early stopping} on a
purged validation block.
\item Fix every source of randomness (initialisation, \omterm{def:ml:neural-networks-for-noisy-tabular-data:dropout}{dropout}, batch order) to one seed per run; train several seeds of
anything you compare, and report the mean and the spread.
\item Deploy the seed ensemble, not the best seed.
\item Compare with ridge and boosted trees on the same months with the same report.
\end{enumerate}
\end{method}

\section{Networks against boosting on tabular data}

The comparison with chapter~4 is on equal terms: same panel, same test years, same report. A single small network
(IC 0.060, net Sharpe ratio 1.37--1.78 across seeds) is ridge regression's equal (0.060, 1.59); its ten-seed ensemble
(0.066, 1.73) matches the lasso's IC and trails the firm's default boosted trees (0.075, 2.02). Gu, Kelly and Xiu found
networks slightly ahead of trees on real US stocks with sixty years of data, and Grinsztajn, Oyallon and Varoquaux trees
ahead of networks on most tabular benchmarks: on panels of ranked characteristics the two families are close, the
network needs an ensemble to be competitive, and the tree ensemble is cheaper to tune.

\section{Tutorial: twenty seeds}\label{tut:ml:neural-networks-for-noisy-tabular-data:tut}

\begin{tutorial}
\textbf{Goal.} Train small and larger networks with ten seeds each, ensemble them, try four regularisers, and compare
with ridge and boosted trees. \textbf{End state:} \cref{fig:ml:nets:history,fig:ml:nets:seeds}, the chapter's numbers.
\begin{tutsteps}
\item \textbf{The deterministic training loop}: one generator for the batch order, the seed for initialisation and
\omterm{def:ml:neural-networks-for-noisy-tabular-data:dropout}{dropout}, \omterm{def:ml:trees-and-boosting:early}{early stopping} on the validation loss.
\omcode{code/firm/nettab/firm_nettab.py}{77}{115}{A deterministic training loop with early stopping.}\label{lst:ml:nets:fit}
\item \textbf{Run} \texttt{ml\_nets.seed\_study()}, \texttt{ensemble()}, \texttt{ablation()}, \texttt{baselines()} and
\texttt{fig\_nets.py}.
\end{tutsteps}
\textbf{What to change next.} Train with 30 seeds per architecture and check \cref{prop:ml:nets:seeds}'s prediction;
replace ReLU by tanh and see whether the seed spread changes.
\end{tutorial}

\section{Build: the firm's tabular networks}\label{bld:ml:neural-networks-for-noisy-tabular-data:nettab}

\begin{build}
\bfield{Purpose} Every network the firm trains on tabular data is reproducible bit for bit and compared across seeds.
\bfield{Interface} \texttt{set\_determinism(threads)}, \texttt{MLP(p, hidden, dropout, norm, activation)},
\texttt{fit(X, y, Xv, yv, hidden, dropout, norm, lr, weight\_decay, batch, epochs, patience, seed) -> Fitted} with
\texttt{predict}, \texttt{history}, \texttt{best\_epoch}, \texttt{state\_hash()}; \texttt{SeedEnsemble(seeds, **kw)}.
\bfield{Rules} One thread, deterministic algorithms, every random draw from the run's seed; inputs standardised with
training statistics stored in the model; the validation block purged by the caller.
\bfield{Acceptance tests} \texttt{code/firm/nettab/tests/}: two runs with one seed give identical weight hashes, two seeds
different ones; the network learns a planted nonlinear function; \omterm{def:ml:trees-and-boosting:early}{early stopping} keeps the best \omterm{def:ml:neural-networks-for-noisy-tabular-data:backprop}{epoch}'s weights; the
ensemble averages its members.
\bfield{Stretch} Mixed-precision training and a data loader for larger panels (chapter~23); export of the weights for
chapter~26's quantised inference.
\end{build}

\begin{omsources}
\item D.~P.~Kingma and J.~Ba, ``Adam: a method for stochastic optimization'', \emph{ICLR}, 2015.
\item I.~Loshchilov and F.~Hutter, ``Decoupled weight decay regularization'', \emph{ICLR}, 2019.
\item N.~Srivastava, G.~Hinton, A.~Krizhevsky, I.~Sutskever and R.~Salakhutdinov, ``Dropout: a simple way to prevent
neural networks from overfitting'', \emph{JMLR} 15, 2014.
\item S.~Ioffe and C.~Szegedy, ``Batch normalization'', \emph{ICML}, 2015; J.~L.~Ba, J.~R.~Kiros and G.~E.~Hinton,
``Layer normalization'', arXiv, 2016.
\item B.~Lakshminarayanan, A.~Pritzel and C.~Blundell, ``Simple and scalable predictive uncertainty estimation using deep
ensembles'', \emph{NeurIPS}, 2017.
\item D.~E.~Rumelhart, G.~E.~Hinton and R.~J.~Williams, ``Learning representations by back-propagating errors'',
\emph{Nature} 323, 1986.
\end{omsources}

\section{Exercises}

\begin{exercise}[$\star$]\label{exo:ml:neural-networks-for-noisy-tabular-data:1}
Count the parameters of an MLP with 40 inputs, hidden layers of 32 and 16 units, and one output. And with a first layer
of 64 units added?
\end{exercise}

\begin{exercise}[$\star$]\label{exo:ml:neural-networks-for-noisy-tabular-data:2}
Two architectures' test ICs vary across seeds with a standard deviation of 0.003 and differ by 0.001 on average. How many
seeds per architecture does a two-standard-error comparison need?
\end{exercise}

\begin{exercise}[$\star$]\label{exo:ml:neural-networks-for-noisy-tabular-data:3}
With 160\,000 training rows and \omterm{def:ml:neural-networks-for-noisy-tabular-data:backprop}{mini-batches} of 512, how many gradient steps is one \omterm{def:ml:neural-networks-for-noisy-tabular-data:backprop}{epoch}? And the chapter's
100\,000 rows?
\end{exercise}

\begin{exercise}[$\star\star$]\label{exo:ml:neural-networks-for-noisy-tabular-data:4}
Why does the seed move R-squared (0.08\% to 0.30\%) so much more than the IC (0.057 to 0.065)?
\end{exercise}

\begin{exercise}[$\star\star$]\label{exo:ml:neural-networks-for-noisy-tabular-data:5}
Why can \omterm{def:ml:neural-networks-for-noisy-tabular-data:dropout}{batch normalisation} hurt a shallow network on ranked inputs, and what would you expect for a deep one?
\end{exercise}

\begin{exercise}[$\star\star$]\label{exo:ml:neural-networks-for-noisy-tabular-data:6}
\emph{Find the flaw.} ``We trained twenty seeds and deployed the one with the best validation score.''
\end{exercise}

\begin{exercise}[$\star\star\star$]\label{exo:ml:neural-networks-for-noisy-tabular-data:7}
\emph{Coding.} Compute the IC of ensembles of the first 2, 5 and 10 seeds of the small network. How much of the gain from
one to ten seeds do two seeds already buy?
\end{exercise}

\begin{exercise}[$\star\star\star$]\label{exo:ml:neural-networks-for-noisy-tabular-data:8}
Show that for one ReLU layer and squared loss, \omterm{def:ml:neural-networks-for-noisy-tabular-data:backprop}{backpropagation}'s gradient with respect to $A_1$ is $\sum_i
(\partial\ell_i/\partial z_{1,i})\odot\mathbf 1\{A_1x_i + c_1 > 0\}\,x_i^\top$, and count its cost against the forward
pass.
\end{exercise}

\section{Problem: Twenty Seeds}

\begin{problem}[{Weekend problem --- a network as a random variable}]\label{pb:ml:neural-networks-for-noisy-tabular-data:1}
Chapter~4's panel; networks trained on months 0--199, stopped on months 201--239, scored on months 240--359.

\medskip\noindent\textbf{Part I --- Training.}
\begin{enumerate}
\item How many parameters do the two networks have?
\item After how many \omterm{def:ml:neural-networks-for-noisy-tabular-data:backprop}{epochs} does \omterm{def:ml:trees-and-boosting:early}{early stopping} stop them, and why so early?
\item What do \omterm{def:ml:neural-networks-for-noisy-tabular-data:dropout}{dropout}, \omterm{def:ml:neural-networks-for-noisy-tabular-data:adam}{weight decay} and the two normalisations do to the IC?
\item Why are three seeds per variant too few to rank them?
\end{enumerate}

\medskip\noindent\textbf{Part II --- Seeds.}
\begin{enumerate}[resume]
\item What range of IC, R-squared and net Sharpe ratio do ten seeds of the small network cover?
\item What are the two architectures' mean ICs, the gap, and the pooled seed standard deviation?
\item How many seeds per architecture does a two-standard-error comparison need?
\item What is the standard error of a comparison of one run of each?
\end{enumerate}

\medskip\noindent\textbf{Part III --- Ensembles.}
\begin{enumerate}[resume]
\item What do the two ten-seed ensembles score?
\item Why does the ensemble beat every member?
\item What happens to the architecture question once each network is ensembled?
\item What does an ensemble give beyond a better forecast?
\end{enumerate}

\medskip\noindent\textbf{Part IV --- The verdict.}
\begin{enumerate}[resume]
\item How do the networks compare with ridge, the lasso and boosted trees on the same months?
\item State the \emph{named result}: the seed standard deviation of the test IC against the architecture gap, and the
number of seeds needed to resolve it at two standard errors.
\item Which model would you deploy, and in what form?
\item What should a research log record about a network?
\item How would the picture change with sixty years of data?
\item What does the state hash guarantee, and what does it not?
\item Where does the randomness of a network's training come from?
\item In one sentence: what is a network trained once on noisy data?
\end{enumerate}
\end{problem}

\section{Interview questions}

\begin{interviewq}[$\star$ \iqroles{mle}]\label{iq:ml:neural-networks-for-noisy-tabular-data:1}
What does \omterm{def:ml:neural-networks-for-noisy-tabular-data:backprop}{backpropagation} compute, and why is it cheap?
\end{interviewq}

\begin{interviewq}[$\star$ \iqroles{mle, researcher}]\label{iq:ml:neural-networks-for-noisy-tabular-data:2}
Explain the difference between \omterm{def:ml:neural-networks-for-noisy-tabular-data:adam}{Adam} and AdamW.
\end{interviewq}

\begin{interviewq}[$\star\star$ \iqroles{researcher, mle}]\label{iq:ml:neural-networks-for-noisy-tabular-data:3}
Your new architecture beats the old one by 0.003 of IC in one run each. What do you do before believing it?
\end{interviewq}

\begin{interviewq}[$\star\star$ \iqroles{mle}]\label{iq:ml:neural-networks-for-noisy-tabular-data:4}
How do you make a PyTorch training run reproducible on CPU? What can still break reproducibility?
\end{interviewq}

\begin{interviewq}[$\star\star$ \iqroles{researcher}]\label{iq:ml:neural-networks-for-noisy-tabular-data:5}
\omterm{def:ml:neural-networks-for-noisy-tabular-data:mlp}{Neural network} or \omterm{def:ml:trees-and-boosting:boosting}{gradient boosting} for a cross-sectional return model? Argue both sides.
\end{interviewq}

\begin{interviewq}[$\star\star\star$ \iqroles{researcher, mle}]\label{iq:ml:neural-networks-for-noisy-tabular-data:6}
Why does averaging networks trained from different seeds improve the forecast, and when does it stop helping?
\end{interviewq}
